% =============================================================================
% model.tex — A simple model of delivery choice and the value of physician time.
% % =============================================================================
% model.tex — A simple model of delivery choice and the value of physician time.
% % =============================================================================
% model.tex — A simple model of delivery choice and the value of physician time.
% % =============================================================================
% model.tex — A simple model of delivery choice and the value of physician time.
% \input{model} from appendix.tex. Organizes the mechanisms and the five
% predictions (P1-P5) that the empirical sections confront.
% =============================================================================

\section{A simple model of delivery choice and the value of physician time}
\label{sec:model}
\setcounter{equation}{0}
\renewcommand{\theequation}{A\arabic{equation}}

\subsection{Environment}

Consider a physician attending a pregnancy for which spontaneous labor begins at an uncertain time $\tilde t$. The physician can provide a vaginal delivery, requiring a random duration $\tau_v$ beginning at $\tilde t$, or schedule a cesarean at a chosen time $t$, taking a relatively predictable duration $\tau_c$. The physician receives billed expected compensation $f_v$ for vaginal delivery, including expected paid labor assistance, and $f_c$ for cesarean delivery.

Let $w(t)$ be the physician's shadow cost of a unit of time, lower in regular working hours than at nights, weekends, and holidays. The uncertainty surrounding vaginal delivery may also require reserving time or remaining reachable, at an opportunity cost $\kappa\geq0$. Let $\theta$ be the patient's net clinical benefit of vaginal rather than cesarean delivery, with distribution $F$ across pregnancies. The physician internalizes a fraction $\alpha\in(0,1]$ of this clinical benefit. A negative $\theta$ denotes a pregnancy for which cesarean delivery is clinically preferred.

\subsection{Choice}

Under the simplifying assumptions that labor duration and onset timing are independent and that a planned cesarean can be scheduled during regular working hours, the private convenience premium from scheduling is
\begin{equation}
\Pi = E[w(\tilde t)]\,E[\tau_v]+\kappa-w_0\tau_c,
\label{eq:premium}
\end{equation}
where $w_0$ is the regular-hours opportunity cost. A planned cesarean is chosen when
\begin{equation}
(f_c-f_v)+\Pi>\alpha\theta.
\label{eq:choice}
\end{equation}

With threshold $\theta^*=(f_c-f_v+\Pi)/\alpha$, the stylized cesarean share is $F(\theta^*)$. This representation separates the net billed fee gap from the cost of unschedulable time. The premium can be positive even where $f_c-f_v<0$, provided vaginal attendance is longer, less predictable, or involves substantial on-call costs. The model is intentionally partial-equilibrium: it does not represent hospital capacity, patient bargaining over dates, insurance contract heterogeneity, or the clinical distinction between planned and urgent intrapartum cesareans.

The response to the fee differential is
\begin{equation}
\frac{\partial F(\theta^*)}{\partial(f_c-f_v)}=\frac{f_F(\theta^*)}{\alpha},
\label{eq:feeresponse}
\end{equation}
where $f_F$ is the density of clinical benefit at the threshold. A large convenience premium does not, by itself, imply an economically negligible fee elasticity. Such a conclusion also requires information about the number of pregnancies close to the threshold. Section~\ref{sec:notprice} therefore uses an explicit external calibration; its fee regressions cannot establish price irrelevance.

\subsection{Mechanisms and alternative implications}

The model predicts that planned cesareans will be concentrated at times when $w(t)$ is low and that a sector in which physicians personally bear greater scheduling costs may have steeper calendar gradients. The model also suggests that arranging cesareans before the rising spontaneous-labor hazard can shift deliveries toward earlier gestational ages. An organizational extension in which hospitals rather than individual clinicians face time-dependent costs yields similar patterns. Hence the weekday gradient, the larger for-profit increment, and attenuation at large obstetric services all have multiple potential interpretations.

A pure physician-leisure story has additional, falsifiable implications: longer uninterrupted breaks should produce stronger displacement than isolated rest days, and procedures should accumulate ahead of extended breaks. The lack of those patterns in Section~\ref{sec:capacity} limits that story. A distinct hospital demand-smoothing story predicts the ability to move volume across weeks; the available evidence does not show it. These findings restrict the mechanisms consistent with the observed calendar patterns.

The model is a device for organizing mechanisms, not an identified structural model. Its primitives, including $w(t)$, $\kappa$, $\alpha$, and $F$, are not estimated from the administrative data.

\subsection{Empirical predictions and corresponding evidence}

P1. Fee responsiveness. A negative mean fee gap is compatible with cesarean delivery when the time-cost premium is sufficiently large. The model alone places no upper bound on the fee response. Section~\ref{sec:notprice} therefore compares the observed fee gaps with externally calibrated responses and presents the unstable fee associations in Table~\ref{tab:fees} and Supplementary Appendix Tables~\ref{tab:feegap_ci}--\ref{tab:fewcluster}. The unimplemented 2015 court order supplies no experimental variation.

P2. Calendar organization. Scheduling lowers the cost of operating in regular hours and can generate fewer prelabor cesarean births on weekends and holidays, with observed delivery composition adjusting in the other direction. Figure~\ref{fig:calendar} and Tables~\ref{tab:main_gradient}--\ref{tab:prelabor_lowrisk} provide the day-of-week, business-hour, and timing comparisons. They cannot isolate hospital scheduling from an individual obstetrician's availability preferences.

P3. Capacity and substitutability. With fewer substitutes to cover unpredictable labor, the effective cost of on-call availability may be higher. \satab{tab:heterogeneity} relates the weekend dip to obstetrician density, while the establishment-level Supplementary Appendix Tables~\ref{tab:org_capacity}--\ref{tab:org_capacity_valid} describe associations with delivery volume and beds. Volume and beds entered together are jointly significant under municipality-by-date effects, without a separately significant coefficient; one-at-a-time specifications are stronger. Registered obstetrician hours do not show the same pattern. None of these regressions identifies a causal capacity effect.

P4. Gestational timing. To complete surgery before spontaneous labor, scheduled deliveries may occur at earlier gestational ages. Section~\ref{sec:cost} and \safig{fig:gestation} document the earlier for-profit distribution, and Table~\ref{tab:health} reports its adjusted early-term association. Births in both sectors peak at week 39. This pattern does not identify how any woman's gestational age would change under a different care arrangement.

P5. Maternal preferences. Mothers may request convenient dates and jointly plan delivery with physicians; a stronger education gradient could be consistent with such a channel. \satab{tab:heterogeneity} reports the education interaction. That proxy neither isolates maternal preferences nor rules them out, even when the pattern is more pronounced among mothers with less schooling.
 from appendix.tex. Organizes the mechanisms and the five
% predictions (P1-P5) that the empirical sections confront.
% =============================================================================

\section{A simple model of delivery choice and the value of physician time}
\label{sec:model}
\setcounter{equation}{0}
\renewcommand{\theequation}{A\arabic{equation}}

\subsection{Environment}

Consider a physician attending a pregnancy for which spontaneous labor begins at an uncertain time $\tilde t$. The physician can provide a vaginal delivery, requiring a random duration $\tau_v$ beginning at $\tilde t$, or schedule a cesarean at a chosen time $t$, taking a relatively predictable duration $\tau_c$. The physician receives billed expected compensation $f_v$ for vaginal delivery, including expected paid labor assistance, and $f_c$ for cesarean delivery.

Let $w(t)$ be the physician's shadow cost of a unit of time, lower in regular working hours than at nights, weekends, and holidays. The uncertainty surrounding vaginal delivery may also require reserving time or remaining reachable, at an opportunity cost $\kappa\geq0$. Let $\theta$ be the patient's net clinical benefit of vaginal rather than cesarean delivery, with distribution $F$ across pregnancies. The physician internalizes a fraction $\alpha\in(0,1]$ of this clinical benefit. A negative $\theta$ denotes a pregnancy for which cesarean delivery is clinically preferred.

\subsection{Choice}

Under the simplifying assumptions that labor duration and onset timing are independent and that a planned cesarean can be scheduled during regular working hours, the private convenience premium from scheduling is
\begin{equation}
\Pi = E[w(\tilde t)]\,E[\tau_v]+\kappa-w_0\tau_c,
\label{eq:premium}
\end{equation}
where $w_0$ is the regular-hours opportunity cost. A planned cesarean is chosen when
\begin{equation}
(f_c-f_v)+\Pi>\alpha\theta.
\label{eq:choice}
\end{equation}

With threshold $\theta^*=(f_c-f_v+\Pi)/\alpha$, the stylized cesarean share is $F(\theta^*)$. This representation separates the net billed fee gap from the cost of unschedulable time. The premium can be positive even where $f_c-f_v<0$, provided vaginal attendance is longer, less predictable, or involves substantial on-call costs. The model is intentionally partial-equilibrium: it does not represent hospital capacity, patient bargaining over dates, insurance contract heterogeneity, or the clinical distinction between planned and urgent intrapartum cesareans.

The response to the fee differential is
\begin{equation}
\frac{\partial F(\theta^*)}{\partial(f_c-f_v)}=\frac{f_F(\theta^*)}{\alpha},
\label{eq:feeresponse}
\end{equation}
where $f_F$ is the density of clinical benefit at the threshold. A large convenience premium does not, by itself, imply an economically negligible fee elasticity. Such a conclusion also requires information about the number of pregnancies close to the threshold. Section~\ref{sec:notprice} therefore uses an explicit external calibration; its fee regressions cannot establish price irrelevance.

\subsection{Mechanisms and alternative implications}

The model predicts that planned cesareans will be concentrated at times when $w(t)$ is low and that a sector in which physicians personally bear greater scheduling costs may have steeper calendar gradients. The model also suggests that arranging cesareans before the rising spontaneous-labor hazard can shift deliveries toward earlier gestational ages. An organizational extension in which hospitals rather than individual clinicians face time-dependent costs yields similar patterns. Hence the weekday gradient, the larger for-profit increment, and attenuation at large obstetric services all have multiple potential interpretations.

A pure physician-leisure story has additional, falsifiable implications: longer uninterrupted breaks should produce stronger displacement than isolated rest days, and procedures should accumulate ahead of extended breaks. The lack of those patterns in Section~\ref{sec:capacity} limits that story. A distinct hospital demand-smoothing story predicts the ability to move volume across weeks; the available evidence does not show it. These findings restrict the mechanisms consistent with the observed calendar patterns.

The model is a device for organizing mechanisms, not an identified structural model. Its primitives, including $w(t)$, $\kappa$, $\alpha$, and $F$, are not estimated from the administrative data.

\subsection{Empirical predictions and corresponding evidence}

P1. Fee responsiveness. A negative mean fee gap is compatible with cesarean delivery when the time-cost premium is sufficiently large. The model alone places no upper bound on the fee response. Section~\ref{sec:notprice} therefore compares the observed fee gaps with externally calibrated responses and presents the unstable fee associations in Table~\ref{tab:fees} and Supplementary Appendix Tables~\ref{tab:feegap_ci}--\ref{tab:fewcluster}. The unimplemented 2015 court order supplies no experimental variation.

P2. Calendar organization. Scheduling lowers the cost of operating in regular hours and can generate fewer prelabor cesarean births on weekends and holidays, with observed delivery composition adjusting in the other direction. Figure~\ref{fig:calendar} and Tables~\ref{tab:main_gradient}--\ref{tab:prelabor_lowrisk} provide the day-of-week, business-hour, and timing comparisons. They cannot isolate hospital scheduling from an individual obstetrician's availability preferences.

P3. Capacity and substitutability. With fewer substitutes to cover unpredictable labor, the effective cost of on-call availability may be higher. \satab{tab:heterogeneity} relates the weekend dip to obstetrician density, while the establishment-level Supplementary Appendix Tables~\ref{tab:org_capacity}--\ref{tab:org_capacity_valid} describe associations with delivery volume and beds. Volume and beds entered together are jointly significant under municipality-by-date effects, without a separately significant coefficient; one-at-a-time specifications are stronger. Registered obstetrician hours do not show the same pattern. None of these regressions identifies a causal capacity effect.

P4. Gestational timing. To complete surgery before spontaneous labor, scheduled deliveries may occur at earlier gestational ages. Section~\ref{sec:cost} and \safig{fig:gestation} document the earlier for-profit distribution, and Table~\ref{tab:health} reports its adjusted early-term association. Births in both sectors peak at week 39. This pattern does not identify how any woman's gestational age would change under a different care arrangement.

P5. Maternal preferences. Mothers may request convenient dates and jointly plan delivery with physicians; a stronger education gradient could be consistent with such a channel. \satab{tab:heterogeneity} reports the education interaction. That proxy neither isolates maternal preferences nor rules them out, even when the pattern is more pronounced among mothers with less schooling.
 from appendix.tex. Organizes the mechanisms and the five
% predictions (P1-P5) that the empirical sections confront.
% =============================================================================

\section{A simple model of delivery choice and the value of physician time}
\label{sec:model}
\setcounter{equation}{0}
\renewcommand{\theequation}{A\arabic{equation}}

\subsection{Environment}

Consider a physician attending a pregnancy for which spontaneous labor begins at an uncertain time $\tilde t$. The physician can provide a vaginal delivery, requiring a random duration $\tau_v$ beginning at $\tilde t$, or schedule a cesarean at a chosen time $t$, taking a relatively predictable duration $\tau_c$. The physician receives billed expected compensation $f_v$ for vaginal delivery, including expected paid labor assistance, and $f_c$ for cesarean delivery.

Let $w(t)$ be the physician's shadow cost of a unit of time, lower in regular working hours than at nights, weekends, and holidays. The uncertainty surrounding vaginal delivery may also require reserving time or remaining reachable, at an opportunity cost $\kappa\geq0$. Let $\theta$ be the patient's net clinical benefit of vaginal rather than cesarean delivery, with distribution $F$ across pregnancies. The physician internalizes a fraction $\alpha\in(0,1]$ of this clinical benefit. A negative $\theta$ denotes a pregnancy for which cesarean delivery is clinically preferred.

\subsection{Choice}

Under the simplifying assumptions that labor duration and onset timing are independent and that a planned cesarean can be scheduled during regular working hours, the private convenience premium from scheduling is
\begin{equation}
\Pi = E[w(\tilde t)]\,E[\tau_v]+\kappa-w_0\tau_c,
\label{eq:premium}
\end{equation}
where $w_0$ is the regular-hours opportunity cost. A planned cesarean is chosen when
\begin{equation}
(f_c-f_v)+\Pi>\alpha\theta.
\label{eq:choice}
\end{equation}

With threshold $\theta^*=(f_c-f_v+\Pi)/\alpha$, the stylized cesarean share is $F(\theta^*)$. This representation separates the net billed fee gap from the cost of unschedulable time. The premium can be positive even where $f_c-f_v<0$, provided vaginal attendance is longer, less predictable, or involves substantial on-call costs. The model is intentionally partial-equilibrium: it does not represent hospital capacity, patient bargaining over dates, insurance contract heterogeneity, or the clinical distinction between planned and urgent intrapartum cesareans.

The response to the fee differential is
\begin{equation}
\frac{\partial F(\theta^*)}{\partial(f_c-f_v)}=\frac{f_F(\theta^*)}{\alpha},
\label{eq:feeresponse}
\end{equation}
where $f_F$ is the density of clinical benefit at the threshold. A large convenience premium does not, by itself, imply an economically negligible fee elasticity. Such a conclusion also requires information about the number of pregnancies close to the threshold. Section~\ref{sec:notprice} therefore uses an explicit external calibration; its fee regressions cannot establish price irrelevance.

\subsection{Mechanisms and alternative implications}

The model predicts that planned cesareans will be concentrated at times when $w(t)$ is low and that a sector in which physicians personally bear greater scheduling costs may have steeper calendar gradients. The model also suggests that arranging cesareans before the rising spontaneous-labor hazard can shift deliveries toward earlier gestational ages. An organizational extension in which hospitals rather than individual clinicians face time-dependent costs yields similar patterns. Hence the weekday gradient, the larger for-profit increment, and attenuation at large obstetric services all have multiple potential interpretations.

A pure physician-leisure story has additional, falsifiable implications: longer uninterrupted breaks should produce stronger displacement than isolated rest days, and procedures should accumulate ahead of extended breaks. The lack of those patterns in Section~\ref{sec:capacity} limits that story. A distinct hospital demand-smoothing story predicts the ability to move volume across weeks; the available evidence does not show it. These findings restrict the mechanisms consistent with the observed calendar patterns.

The model is a device for organizing mechanisms, not an identified structural model. Its primitives, including $w(t)$, $\kappa$, $\alpha$, and $F$, are not estimated from the administrative data.

\subsection{Empirical predictions and corresponding evidence}

P1. Fee responsiveness. A negative mean fee gap is compatible with cesarean delivery when the time-cost premium is sufficiently large. The model alone places no upper bound on the fee response. Section~\ref{sec:notprice} therefore compares the observed fee gaps with externally calibrated responses and presents the unstable fee associations in Table~\ref{tab:fees} and Supplementary Appendix Tables~\ref{tab:feegap_ci}--\ref{tab:fewcluster}. The unimplemented 2015 court order supplies no experimental variation.

P2. Calendar organization. Scheduling lowers the cost of operating in regular hours and can generate fewer prelabor cesarean births on weekends and holidays, with observed delivery composition adjusting in the other direction. Figure~\ref{fig:calendar} and Tables~\ref{tab:main_gradient}--\ref{tab:prelabor_lowrisk} provide the day-of-week, business-hour, and timing comparisons. They cannot isolate hospital scheduling from an individual obstetrician's availability preferences.

P3. Capacity and substitutability. With fewer substitutes to cover unpredictable labor, the effective cost of on-call availability may be higher. \satab{tab:heterogeneity} relates the weekend dip to obstetrician density, while the establishment-level Supplementary Appendix Tables~\ref{tab:org_capacity}--\ref{tab:org_capacity_valid} describe associations with delivery volume and beds. Volume and beds entered together are jointly significant under municipality-by-date effects, without a separately significant coefficient; one-at-a-time specifications are stronger. Registered obstetrician hours do not show the same pattern. None of these regressions identifies a causal capacity effect.

P4. Gestational timing. To complete surgery before spontaneous labor, scheduled deliveries may occur at earlier gestational ages. Section~\ref{sec:cost} and \safig{fig:gestation} document the earlier for-profit distribution, and Table~\ref{tab:health} reports its adjusted early-term association. Births in both sectors peak at week 39. This pattern does not identify how any woman's gestational age would change under a different care arrangement.

P5. Maternal preferences. Mothers may request convenient dates and jointly plan delivery with physicians; a stronger education gradient could be consistent with such a channel. \satab{tab:heterogeneity} reports the education interaction. That proxy neither isolates maternal preferences nor rules them out, even when the pattern is more pronounced among mothers with less schooling.
 from appendix.tex. Organizes the mechanisms and the five
% predictions (P1-P5) that the empirical sections confront.
% =============================================================================

\section{A simple model of delivery choice and the value of physician time}
\label{sec:model}
\setcounter{equation}{0}
\renewcommand{\theequation}{A\arabic{equation}}

\subsection{Environment}

Consider a physician attending a pregnancy for which spontaneous labor begins at an uncertain time $\tilde t$. The physician can provide a vaginal delivery, requiring a random duration $\tau_v$ beginning at $\tilde t$, or schedule a cesarean at a chosen time $t$, taking a relatively predictable duration $\tau_c$. The physician receives billed expected compensation $f_v$ for vaginal delivery, including expected paid labor assistance, and $f_c$ for cesarean delivery.

Let $w(t)$ be the physician's shadow cost of a unit of time, lower in regular working hours than at nights, weekends, and holidays. The uncertainty surrounding vaginal delivery may also require reserving time or remaining reachable, at an opportunity cost $\kappa\geq0$. Let $\theta$ be the patient's net clinical benefit of vaginal rather than cesarean delivery, with distribution $F$ across pregnancies. The physician internalizes a fraction $\alpha\in(0,1]$ of this clinical benefit. A negative $\theta$ denotes a pregnancy for which cesarean delivery is clinically preferred.

\subsection{Choice}

Under the simplifying assumptions that labor duration and onset timing are independent and that a planned cesarean can be scheduled during regular working hours, the private convenience premium from scheduling is
\begin{equation}
\Pi = E[w(\tilde t)]\,E[\tau_v]+\kappa-w_0\tau_c,
\label{eq:premium}
\end{equation}
where $w_0$ is the regular-hours opportunity cost. A planned cesarean is chosen when
\begin{equation}
(f_c-f_v)+\Pi>\alpha\theta.
\label{eq:choice}
\end{equation}

With threshold $\theta^*=(f_c-f_v+\Pi)/\alpha$, the stylized cesarean share is $F(\theta^*)$. This representation separates the net billed fee gap from the cost of unschedulable time. The premium can be positive even where $f_c-f_v<0$, provided vaginal attendance is longer, less predictable, or involves substantial on-call costs. The model is intentionally partial-equilibrium: it does not represent hospital capacity, patient bargaining over dates, insurance contract heterogeneity, or the clinical distinction between planned and urgent intrapartum cesareans.

The response to the fee differential is
\begin{equation}
\frac{\partial F(\theta^*)}{\partial(f_c-f_v)}=\frac{f_F(\theta^*)}{\alpha},
\label{eq:feeresponse}
\end{equation}
where $f_F$ is the density of clinical benefit at the threshold. A large convenience premium does not, by itself, imply an economically negligible fee elasticity. Such a conclusion also requires information about the number of pregnancies close to the threshold. Section~\ref{sec:notprice} therefore uses an explicit external calibration; its fee regressions cannot establish price irrelevance.

\subsection{Mechanisms and alternative implications}

The model predicts that planned cesareans will be concentrated at times when $w(t)$ is low and that a sector in which physicians personally bear greater scheduling costs may have steeper calendar gradients. The model also suggests that arranging cesareans before the rising spontaneous-labor hazard can shift deliveries toward earlier gestational ages. An organizational extension in which hospitals rather than individual clinicians face time-dependent costs yields similar patterns. Hence the weekday gradient, the larger for-profit increment, and attenuation at large obstetric services all have multiple potential interpretations.

A pure physician-leisure story has additional, falsifiable implications: longer uninterrupted breaks should produce stronger displacement than isolated rest days, and procedures should accumulate ahead of extended breaks. The lack of those patterns in Section~\ref{sec:capacity} limits that story. A distinct hospital demand-smoothing story predicts the ability to move volume across weeks; the available evidence does not show it. These findings restrict the mechanisms consistent with the observed calendar patterns.

The model is a device for organizing mechanisms, not an identified structural model. Its primitives, including $w(t)$, $\kappa$, $\alpha$, and $F$, are not estimated from the administrative data.

\subsection{Empirical predictions and corresponding evidence}

P1. Fee responsiveness. A negative mean fee gap is compatible with cesarean delivery when the time-cost premium is sufficiently large. The model alone places no upper bound on the fee response. Section~\ref{sec:notprice} therefore compares the observed fee gaps with externally calibrated responses and presents the unstable fee associations in Table~\ref{tab:fees} and Supplementary Appendix Tables~\ref{tab:feegap_ci}--\ref{tab:fewcluster}. The unimplemented 2015 court order supplies no experimental variation.

P2. Calendar organization. Scheduling lowers the cost of operating in regular hours and can generate fewer prelabor cesarean births on weekends and holidays, with observed delivery composition adjusting in the other direction. Figure~\ref{fig:calendar} and Tables~\ref{tab:main_gradient}--\ref{tab:prelabor_lowrisk} provide the day-of-week, business-hour, and timing comparisons. They cannot isolate hospital scheduling from an individual obstetrician's availability preferences.

P3. Capacity and substitutability. With fewer substitutes to cover unpredictable labor, the effective cost of on-call availability may be higher. \satab{tab:heterogeneity} relates the weekend dip to obstetrician density, while the establishment-level Supplementary Appendix Tables~\ref{tab:org_capacity}--\ref{tab:org_capacity_valid} describe associations with delivery volume and beds. Volume and beds entered together are jointly significant under municipality-by-date effects, without a separately significant coefficient; one-at-a-time specifications are stronger. Registered obstetrician hours do not show the same pattern. None of these regressions identifies a causal capacity effect.

P4. Gestational timing. To complete surgery before spontaneous labor, scheduled deliveries may occur at earlier gestational ages. Section~\ref{sec:cost} and \safig{fig:gestation} document the earlier for-profit distribution, and Table~\ref{tab:health} reports its adjusted early-term association. Births in both sectors peak at week 39. This pattern does not identify how any woman's gestational age would change under a different care arrangement.

P5. Maternal preferences. Mothers may request convenient dates and jointly plan delivery with physicians; a stronger education gradient could be consistent with such a channel. \satab{tab:heterogeneity} reports the education interaction. That proxy neither isolates maternal preferences nor rules them out, even when the pattern is more pronounced among mothers with less schooling.
